\documentclass[12pt,a4paper]{article}
\usepackage[UTF8]{ctex}
\usepackage{amsmath,amssymb,amsthm}
\usepackage{geometry}
\usepackage{graphicx}
\usepackage{hyperref}
\usepackage{bm}

\geometry{margin=2.5cm}

\title{Body雅可比与空间雅可比详解}
\author{Modern Robotics}
\date{}

\begin{document}

\maketitle

\tableofcontents
\newpage

\section{引言}

在机器人运动学和动力学中，雅可比矩阵（Jacobian）是一个核心概念，它将关节速度映射到末端执行器（或连杆）的速度。根据参考坐标系的不同，雅可比矩阵有两种主要形式：

\begin{itemize}
\item \textbf{Body雅可比}（Body Jacobian）：在运动坐标系（体坐标系）中表示
\item \textbf{空间雅可比}（Space Jacobian）：在固定坐标系（空间坐标系）中表示
\end{itemize}

本文将详细介绍这两种雅可比矩阵的定义、构造方法、相互关系及其在机器人学中的应用。

\section{雅可比矩阵的基本概念}

\subsection{定义}

对于$n$关节机器人，雅可比矩阵$J(\theta) \in \mathbb{R}^{6 \times n}$将关节速度$\dot{\theta} \in \mathbb{R}^n$映射到末端执行器的twist $V \in \mathbb{R}^6$：

\begin{equation}
V = J(\theta) \dot{\theta}
\label{eq:jacobian_basic}
\end{equation}

其中twist $V = \begin{bmatrix} \omega \\ v \end{bmatrix}$包含角速度$\omega \in \mathbb{R}^3$和线速度$v \in \mathbb{R}^3$。

\subsection{物理意义}

雅可比矩阵的每一列对应一个关节对末端执行器速度的贡献：

\begin{itemize}
\item 第$j$列$J_j(\theta)$表示：当只有关节$j$以单位速度运动时，末端执行器产生的twist
\item 第$j$列就是关节$j$的螺旋轴（screw axis）在相应坐标系中的表示
\end{itemize}

\subsection{两种表示形式}

根据twist $V$在哪个坐标系中表示，雅可比矩阵有两种形式：

\begin{enumerate}
\item \textbf{Body雅可比} $J^b(\theta)$：twist在末端执行器坐标系（体坐标系）中表示
\item \textbf{空间雅可比} $J^s(\theta)$：twist在基座坐标系（空间坐标系）中表示
\end{enumerate}

\section{Body雅可比（Body Jacobian）}

\subsection{定义}

\textbf{Body雅可比} $J_i^b(\theta)$将关节速度$\dot{\theta}$映射到连杆$i$的Body Twist $V_i^b$：

\begin{equation}
V_i^b = J_i^b(\theta) \dot{\theta}
\label{eq:body_jacobian_def}
\end{equation}

其中$V_i^b$是连杆$i$在坐标系$\{i\}$中表示的twist。

\subsection{从$T_{0i}^{-1}\dot{T}_{0i}$提取Body雅可比}

设$T_{0i}(\theta_1, \ldots, \theta_i)$表示从基座坐标系$\{0\}$到连杆坐标系$\{i\}$的正向运动学变换。

\textbf{关键关系}：从$T_{0i}^{-1}\dot{T}_{0i}$可以提取Body Twist：

\begin{equation}
T_{0i}^{-1} \dot{T}_{0i} = [V_i^b]
\label{eq:T_inv_T_dot}
\end{equation}

其中$[V_i^b]$是Body Twist的反对称矩阵表示。

\subsection{Body雅可比的构造}

\textbf{核心结论}：Body雅可比$J_i^b(\theta)$的第$j$列是$A_j^b$（$j = 1, \ldots, i$），其余列为零。

\textbf{推导}：

对$T_{0i}(\theta)$关于$\theta_j$求偏导（$j = 1, \ldots, i$）：

\begin{equation}
\frac{\partial T_{0i}}{\partial \theta_j} = T_{0i} [A_j^b]
\label{eq:partial_T}
\end{equation}

其中$A_j^b$是关节$j$的螺旋轴，在坐标系$\{i\}$中表示。

因此：
\begin{equation}
T_{0i}^{-1} \frac{\partial T_{0i}}{\partial \theta_j} = [A_j^b]
\end{equation}

\textbf{从$\dot{T}_{0i}$展开}：

\begin{align}
\dot{T}_{0i} &= \sum_{j=1}^{i} \frac{\partial T_{0i}}{\partial \theta_j} \dot{\theta}_j \\
&= \sum_{j=1}^{i} T_{0i} [A_j^b] \dot{\theta}_j
\end{align}

因此：
\begin{align}
T_{0i}^{-1} \dot{T}_{0i} &= \sum_{j=1}^{i} [A_j^b] \dot{\theta}_j \\
&= \left[ \sum_{j=1}^{i} A_j^b \dot{\theta}_j \right] \\
&= [V_i^b]
\end{align}

\textbf{因此}：
\begin{equation}
V_i^b = \sum_{j=1}^{i} A_j^b \dot{\theta}_j = J_i^b(\theta) \dot{\theta}
\end{equation}

其中Body雅可比$J_i^b(\theta)$的第$j$列是$A_j^b$（$j = 1, \ldots, i$），其余列为零。

\subsection{完整Body雅可比的构造}

由于$J_i^b$需要是$6 \times n$矩阵（$n$是总关节数），而$T_{0i}$只依赖于前$i$个关节，我们：

\begin{enumerate}
\item 计算前$i$列：从$T_{0i}^{-1}\dot{T}_{0i}$提取，每列对应一个关节的螺旋轴$A_j^b$
\item 填充后$n-i$列：全部为零（因为后$n-i$个关节不影响连杆$i$的运动）
\end{enumerate}

\textbf{因此}：
\begin{equation}
J_i^b(\theta) = \begin{bmatrix} A_1^b & A_2^b & \cdots & A_i^b & 0 & \cdots & 0 \end{bmatrix} \in \mathbb{R}^{6 \times n}
\label{eq:body_jacobian_full}
\end{equation}

\subsection{Body雅可比的计算方法}

\textbf{方法1：从偏导数计算}

\begin{enumerate}
\item 计算正向运动学$T_{0i}(\theta)$
\item 对每个关节角度$\theta_j$（$j = 1, \ldots, i$）求偏导
\item 计算$T_{0i}^{-1} \frac{\partial T_{0i}}{\partial \theta_j}$，提取螺旋轴$A_j^b$
\item 构造$J_i^b$的第$j$列为$A_j^b$
\item 后$n-i$列填充为零
\end{enumerate}

\textbf{方法2：从$T_{0i}^{-1}\dot{T}_{0i}$提取}

\begin{enumerate}
\item 计算$T_{0i}(\theta)$
\item 计算$\dot{T}_{0i} = \sum_{j=1}^{i} \frac{\partial T_{0i}}{\partial \theta_j} \dot{\theta}_j$
\item 计算$T_{0i}^{-1} \dot{T}_{0i} = [V_i^b]$
\item 从$[V_i^b]$中提取$V_i^b = \sum_{j=1}^{i} A_j^b \dot{\theta}_j$
\item 识别出各列对应的$A_j^b$
\end{enumerate}

\section{空间雅可比（Space Jacobian）}

\subsection{定义}

\textbf{空间雅可比} $J^s(\theta)$将关节速度$\dot{\theta}$映射到末端执行器的空间Twist $V^s$：

\begin{equation}
V^s = J^s(\theta) \dot{\theta}
\label{eq:space_jacobian_def}
\end{equation}

其中$V^s$是末端执行器在基座坐标系$\{0\}$中表示的twist。

\subsection{从$\dot{T} = [V^s]T$提取空间雅可比}

对于末端执行器的变换矩阵$T_{0n}(\theta)$，有：

\begin{equation}
\dot{T}_{0n} = [V^s] T_{0n}
\label{eq:T_dot_space}
\end{equation}

其中$[V^s]$是空间Twist的反对称矩阵表示。

\subsection{空间雅可比的构造}

\textbf{核心结论}：空间雅可比$J^s(\theta)$的第$j$列是$A_j^s$（$j = 1, \ldots, n$），其中$A_j^s$是关节$j$的螺旋轴在基座坐标系$\{0\}$中的表示。

\textbf{推导}：

对$T_{0n}(\theta)$关于$\theta_j$求偏导：

\begin{equation}
\frac{\partial T_{0n}}{\partial \theta_j} = [A_j^s] T_{0n}
\label{eq:partial_T_space}
\end{equation}

其中$A_j^s$是关节$j$的螺旋轴，在基座坐标系$\{0\}$中表示。

\textbf{从$\dot{T}_{0n}$展开}：

\begin{align}
\dot{T}_{0n} &= \sum_{j=1}^{n} \frac{\partial T_{0n}}{\partial \theta_j} \dot{\theta}_j \\
&= \sum_{j=1}^{n} [A_j^s] T_{0n} \dot{\theta}_j
\end{align}

因此：
\begin{align}
\dot{T}_{0n} T_{0n}^{-1} &= \sum_{j=1}^{n} [A_j^s] \dot{\theta}_j \\
&= \left[ \sum_{j=1}^{n} A_j^s \dot{\theta}_j \right] \\
&= [V^s]
\end{align}

\textbf{因此}：
\begin{equation}
V^s = \sum_{j=1}^{n} A_j^s \dot{\theta}_j = J^s(\theta) \dot{\theta}
\end{equation}

其中空间雅可比$J^s(\theta)$的第$j$列是$A_j^s$。

\subsection{完整空间雅可比的构造}

\begin{equation}
J^s(\theta) = \begin{bmatrix} A_1^s & A_2^s & \cdots & A_n^s \end{bmatrix} \in \mathbb{R}^{6 \times n}
\label{eq:space_jacobian_full}
\end{equation}

其中$A_j^s$是关节$j$的螺旋轴在基座坐标系$\{0\}$中的表示。

\subsection{空间雅可比的计算方法}

\textbf{方法1：从偏导数计算}

\begin{enumerate}
\item 计算正向运动学$T_{0n}(\theta)$
\item 对每个关节角度$\theta_j$（$j = 1, \ldots, n$）求偏导
\item 计算$\frac{\partial T_{0n}}{\partial \theta_j} T_{0n}^{-1}$，提取螺旋轴$A_j^s$
\item 构造$J^s$的第$j$列为$A_j^s$
\end{enumerate}

\textbf{方法2：从$\dot{T}_{0n} T_{0n}^{-1}$提取}

\begin{enumerate}
\item 计算$T_{0n}(\theta)$
\item 计算$\dot{T}_{0n} = \sum_{j=1}^{n} \frac{\partial T_{0n}}{\partial \theta_j} \dot{\theta}_j$
\item 计算$\dot{T}_{0n} T_{0n}^{-1} = [V^s]$
\item 从$[V^s]$中提取$V^s = \sum_{j=1}^{n} A_j^s \dot{\theta}_j$
\item 识别出各列对应的$A_j^s$
\end{enumerate}

\section{Body雅可比与空间雅可比的关系}

\subsection{Twist的转换}

Body Twist $V^b$和空间Twist $V^s$通过伴随变换（adjoint transformation）关联：

\begin{equation}
V^s = \text{Ad}_{T_{0n}}(V^b) = \begin{bmatrix} R & 0 \\ [p]R & R \end{bmatrix} V^b
\label{eq:twist_conversion}
\end{equation}

其中$T_{0n} = \begin{bmatrix} R & p \\ 0 & 1 \end{bmatrix}$，$\text{Ad}_{T_{0n}}$是伴随变换矩阵。

\subsection{雅可比矩阵的转换}

由于$V^s = J^s(\theta) \dot{\theta}$和$V^b = J^b(\theta) \dot{\theta}$，以及$V^s = \text{Ad}_{T_{0n}}(V^b)$，我们有：

\begin{align}
J^s(\theta) \dot{\theta} &= \text{Ad}_{T_{0n}}(J^b(\theta) \dot{\theta}) \\
&= \text{Ad}_{T_{0n}} J^b(\theta) \dot{\theta}
\end{align}

\textbf{因此}：
\begin{equation}
J^s(\theta) = \text{Ad}_{T_{0n}} J^b(\theta)
\label{eq:jacobian_conversion}
\end{equation}

\textbf{逆关系}：
\begin{equation}
J^b(\theta) = \text{Ad}_{T_{0n}}^{-1} J^s(\theta) = \text{Ad}_{T_{0n}^{-1}} J^s(\theta)
\label{eq:jacobian_conversion_inv}
\end{equation}

\subsection{列向量的转换}

对于第$j$列，有：

\begin{equation}
A_j^s = \text{Ad}_{T_{0n}}(A_j^b)
\label{eq:column_conversion}
\end{equation}

这表示：关节$j$的螺旋轴在基座坐标系中的表示等于其在末端执行器坐标系中的表示通过伴随变换得到。

\section{具体例子}

\subsection{例子1：2R平面机器人}

考虑2R平面机器人，两个关节都在$xy$平面内旋转。

\textbf{正向运动学}：
\begin{equation}
T_{02}(\theta_1, \theta_2) = T_{01}(\theta_1) T_{12}(\theta_2)
\end{equation}

\textbf{Body雅可比} $J_2^b(\theta)$：

\begin{itemize}
\item 第1列$A_1^b$：关节1的螺旋轴在坐标系$\{2\}$中的表示
\item 第2列$A_2^b$：关节2的螺旋轴在坐标系$\{2\}$中的表示（通常就是$A_2$本身）
\end{itemize}

\textbf{空间雅可比} $J^s(\theta)$：

\begin{itemize}
\item 第1列$A_1^s$：关节1的螺旋轴在基座坐标系$\{0\}$中的表示
\item 第2列$A_2^s$：关节2的螺旋轴在基座坐标系$\{0\}$中的表示
\end{itemize}

\textbf{关系}：
\begin{equation}
J^s(\theta) = \text{Ad}_{T_{02}} J_2^b(\theta)
\end{equation}

\subsection{例子2：单关节旋转}

考虑绕$z$轴旋转的单个关节，旋转角度为$\theta$。

\textbf{Body雅可比}：
\begin{equation}
J^b(\theta) = A_1^b = \begin{bmatrix} 0 \\ 0 \\ 1 \\ 0 \\ 0 \\ 0 \end{bmatrix}
\end{equation}

（绕$z$轴的螺旋轴在体坐标系中表示）

\textbf{空间雅可比}：
\begin{equation}
J^s(\theta) = A_1^s = \begin{bmatrix} 0 \\ 0 \\ 1 \\ 0 \\ 0 \\ 0 \end{bmatrix}
\end{equation}

（如果基座坐标系的$z$轴与关节轴对齐，则两者相同）

\section{在机器人动力学中的应用}

\subsection{动能表达式}

机器人的总动能可以表示为：

\begin{equation}
K = \frac{1}{2}\sum_{i=1}^{n} (V_i^b)^T G_i V_i^b
\end{equation}

其中$G_i$是连杆$i$的空间惯性矩阵（在坐标系$\{i\}$中表示）。

使用$V_i^b = J_i^b(\theta) \dot{\theta}$：

\begin{equation}
K = \frac{1}{2}\dot{\theta}^T \left(\sum_{i=1}^{n} J_i^b(\theta)^T G_i J_i^b(\theta)\right) \dot{\theta} = \frac{1}{2}\dot{\theta}^T M(\theta) \dot{\theta}
\end{equation}

其中质量矩阵为：
\begin{equation}
M(\theta) = \sum_{i=1}^{n} J_i^b(\theta)^T G_i J_i^b(\theta)
\label{eq:mass_matrix}
\end{equation}

\subsection{力映射}

雅可比矩阵的转置用于将任务空间的wrench映射到关节空间的力矩：

\begin{equation}
\tau = J^T(\theta) F
\label{eq:force_mapping}
\end{equation}

其中：
\begin{itemize}
\item $F \in \mathbb{R}^6$：作用在末端执行器上的wrench（与$V$在同一坐标系中表示）
\item $\tau \in \mathbb{R}^n$：关节力矩
\item $J(\theta)$：可以是Body雅可比或空间雅可比（取决于$F$的表示）
\end{itemize}

\subsection{速度控制}

在速度控制中，使用雅可比矩阵的逆或伪逆：

\begin{equation}
\dot{\theta} = J^{-1}(\theta) V \quad \text{（当$J$可逆时）}
\end{equation}

或

\begin{equation}
\dot{\theta} = J^{\dagger}(\theta) V \quad \text{（使用伪逆）}
\end{equation}

其中$J^{\dagger} = (J^T J)^{-1}J^T$是伪逆。

\section{Body雅可比与空间雅可比的比较}

\subsection{相同点}

\begin{itemize}
\item 都是$6 \times n$矩阵
\item 都将关节速度映射到twist
\item 奇异性相同（$\det(J^b) = 0$当且仅当$\det(J^s) = 0$）
\item 秩相同（$\text{rank}(J^b) = \text{rank}(J^s)$）
\end{itemize}

\subsection{不同点}

\begin{table}[h]
\centering
\begin{tabular}{|l|l|l|}
\hline
\textbf{特性} & \textbf{Body雅可比} & \textbf{空间雅可比} \\
\hline
参考坐标系 & 末端执行器坐标系（体坐标系） & 基座坐标系（空间坐标系） \\
\hline
Twist表示 & $V^b$（Body Twist） & $V^s$（空间Twist） \\
\hline
提取方法 & $T_{0i}^{-1}\dot{T}_{0i} = [V_i^b]$ & $\dot{T}_{0n} T_{0n}^{-1} = [V^s]$ \\
\hline
列向量 & $A_j^b$（在坐标系$\{i\}$中） & $A_j^s$（在坐标系$\{0\}$中） \\
\hline
偏导数关系 & $\frac{\partial T_{0i}}{\partial \theta_j} = T_{0i} [A_j^b]$ & $\frac{\partial T_{0n}}{\partial \theta_j} = [A_j^s] T_{0n}$ \\
\hline
适用场景 & 动力学计算、力控制 & 运动学分析、轨迹规划 \\
\hline
\end{tabular}
\caption{Body雅可比与空间雅可比的比较}
\label{tab:comparison}
\end{table}

\subsection{选择建议}

\begin{itemize}
\item \textbf{使用Body雅可比}：
\begin{itemize}
\item 计算质量矩阵$M(\theta)$
\item 动力学分析和控制
\item 力控制（当力在体坐标系中表示时）
\end{itemize}

\item \textbf{使用空间雅可比}：
\begin{itemize}
\item 运动学分析和可视化
\item 轨迹规划（当轨迹在基座坐标系中表示时）
\item 奇异点分析
\end{itemize}
\end{itemize}

\section{总结}

\subsection{关键结论}

\begin{enumerate}
\item \textbf{Body雅可比} $J_i^b(\theta)$：
\begin{itemize}
\item 第$j$列是$A_j^b$（$j = 1, \ldots, i$），其余列为零
\item 从$T_{0i}^{-1}\dot{T}_{0i}$提取
\item 将关节速度映射到Body Twist
\end{itemize}

\item \textbf{空间雅可比} $J^s(\theta)$：
\begin{itemize}
\item 第$j$列是$A_j^s$（$j = 1, \ldots, n$）
\item 从$\dot{T}_{0n} T_{0n}^{-1}$提取
\item 将关节速度映射到空间Twist
\end{itemize}

\item \textbf{关系}：
\begin{itemize}
\item $J^s(\theta) = \text{Ad}_{T_{0n}} J^b(\theta)$
\item $A_j^s = \text{Ad}_{T_{0n}}(A_j^b)$
\end{itemize}
\end{enumerate}

\subsection{物理意义}

\begin{itemize}
\item 雅可比矩阵的每一列对应一个关节对末端执行器速度的贡献
\item Body雅可比在运动坐标系中观察速度
\item 空间雅可比在固定坐标系中观察速度
\item 两者通过伴随变换相互转换
\end{itemize}

\subsection{应用}

\begin{itemize}
\item \textbf{运动学}：速度映射、逆运动学
\item \textbf{动力学}：质量矩阵构造、动能计算
\item \textbf{控制}：速度控制、力控制
\item \textbf{分析}：奇异点分析、可操作性分析
\end{itemize}

\section{参考文献}

\begin{itemize}
\item Modern Robotics: Mechanics, Planning, and Control
\item 第3章：刚体运动（Twist和Screw理论）
\item 第5章：速度运动学和静力学（雅可比矩阵）
\item 第8章：开链动力学
\end{itemize}

\end{document}
